\hypertarget{archive-formats-zipfile-tarfile}{%
\chapter{\texorpdfstring{Archive Formats (\texttt{zipfile},
\texttt{tarfile})}{Archive Formats (zipfile, tarfile)}}\label{archive-formats-zipfile-tarfile}}

Archives group multiple files into a single container, often with
compression.

\hypertarget{zipfile-the-directory-at-the-end-architecture}{%
\subsection{\texorpdfstring{50.1 \texttt{zipfile}: The
Directory-at-the-End
Architecture}{50.1 zipfile: The Directory-at-the-End Architecture}}\label{zipfile-the-directory-at-the-end-architecture}}

The ZIP format stores its \textbf{Central Directory} at the \emph{end}
of the file. * \textbf{Random Access}: This design allows a program to
read the directory once and then jump (seek) to any file within the
archive without reading the whole file. * \textbf{Encryption}:
\passthrough{\lstinline!zipfile!} supports password-protected archives
(Legacy ZIP encryption), but it is computationally weak. For modern
security, use an external library like
\passthrough{\lstinline!pycryptodome!} for AES-256.

\hypertarget{tarfile-the-tape-archive-heritage}{%
\subsection{\texorpdfstring{50.2 \texttt{tarfile}: The Tape Archive
Heritage}{50.2 tarfile: The Tape Archive Heritage}}\label{tarfile-the-tape-archive-heritage}}

Originally designed for magnetic tapes, the TAR format is a simple
concatenation of files with a 512-byte header for each. * \textbf{No
Central Directory}: To find a file at the end of a
\passthrough{\lstinline!.tar!} archive, you must read all previous
headers. * \textbf{Compression}: \passthrough{\lstinline!.tar.gz!} or
\passthrough{\lstinline!.tar.xz!} are created by piping the output of
the TAR stream into a compressor. \passthrough{\lstinline!tarfile!}
handles this transparently via its \passthrough{\lstinline!mode!}
argument (e.g., \passthrough{\lstinline!'w:gz'!}). * \textbf{Sparse
Files}: \passthrough{\lstinline!tarfile!} can handle ``sparse files''
(files with large holes of zeros) efficiently, preserving their
structure on disk without allocating physical space for the zeros.

\begin{center}\rule{0.5\linewidth}{0.5pt}\end{center}
